%───────────────────────────────────────────────────────────────────────────────
% METADATOS DEL DOCUMENTO — CURSO PROGRAMMING LANGUAGES (CIIC4030-066)
%───────────────────────────────────────────────────────────────────────────────

% --- Identificación del documento ---
\newcommand{\docTitulo}{Fundamentos y Mecánica de los Lenguajes de Programación}
\newcommand{\docSubtitulo}{De la teoría de la computación\\ al diseño de sistemas\\ y algoritmos de alta precisión}
\newcommand{\docAutor}{Jesús Antonio Chala Casanova}
\newcommand{\docFecha}{2026} % Basado en la fecha actual
\newcommand{\docEstado}{Notas de Curso (CISE-Ph.D.)}

% --- Cita de portada ---
% Sipser enfatiza que la teoría proporciona las herramientas conceptuales 
% necesarias para el diseño técnico [3].
\newcommand{\docCita}{}

% --- Páginas preliminares ---
\newcommand{\docDedicatoria}{}
\newcommand{\docEpigrafe}{``El diseño de un lenguaje de programación es una búsqueda de abstracciones que permitan al intelecto humano dominar la complejidad del cómputo.''}
\newcommand{\docAgradecimientos}{}

% --- Metadatos PDF ---
\newcommand{\docKeywords}{compilación, autómatas finitos, gramáticas libres de contexto, cálculo lambda, programación funcional, gestión de memoria, Scala, Go, Michael Sipser, computación científica}
\newcommand{\docSubject}{Notas de estudio avanzado sobre lenguajes de programación y teoría de la computación}

% =============================================================================
% RECURSOS BIBLIOGRÁFICOS
% =============================================================================
\addbibresource{references/references.bib}
